\chapter{The headline experiment}
\label{ch:headline}

\section{The answer, first}

\begin{quote}
Does a fifty kilobyte linear model that runs in microseconds on a CPU get close
enough to a twelve billion parameter model to be the right default for a
developer tool?
\end{quote}

\noindent\textbf{One correction to the question before answering it.} The size in
it is the build specification's own framing, written before the model existed,
and the exported graph is larger than that. The artefact is committed under
\code{models/} and anybody can check its size directly. Nothing in the answer
turns on the figure: the argument is about a linear model that runs on a CPU in
microseconds per field against a twelve billion parameter model on an
accelerator, and that contrast is measured rather than asserted.

\textbf{The premise did not survive the measurement. The linear model did not
have to get close, because it beat the language model on both averages. The rule
table beat both.}

\begin{table}[htbp]
  \centering
  \caption{The headline benchmark: three engines on the test split, both
  averages, and the fourth row that does not exist.}
  \label{tab:headline}
  \begingroup\small\setlength{\tabcolsep}{4pt}
\begin{tabular}{lrrrrrr}
\toprule
\textbf{Engine} & \textbf{Fields} & \textbf{Forms} & \textbf{Templates} & \textbf{Labels} & \textbf{macro-F1} & \textbf{micro-F1} \\  % results: experiments/results/bench/2026-08-26T20-32-48Z_bench_6457ac7/benchmark.json
\midrule
\texttt{rules} & 2317 & 240 & 10 & 36 & 0.7911 & 0.7678 \\  % results: experiments/results/bench/2026-08-26T20-32-48Z_bench_6457ac7/benchmark.json
\texttt{ngram} & 2317 & 240 & 10 & 39 & 0.7394 & 0.8226 \\  % results: experiments/results/bench/2026-08-26T20-32-48Z_bench_6457ac7/benchmark.json
\texttt{llm} & 2317 & 240 & 10 & 38 & 0.6353 & 0.6953 \\  % results: experiments/results/bench/2026-08-26T20-32-48Z_bench_6457ac7/benchmark.json
\midrule
\texttt{bert-onnx-int8} & absent & absent & absent & absent & absent & absent \\  % results: experiments/results/bench/2026-08-26T20-32-48Z_bench_6457ac7/benchmark.json
\bottomrule
\end{tabular}
\endgroup
\par\smallskip\noindent{\footnotesize Source: \rpath{experiments/results/bench/2026-08-26T20-32-48Z_bench_6457ac7/benchmark.json}}

\end{table}

\begin{figure}[htbp]
  \centering
  \includegraphics[width=.78\textwidth]{engine_f1.pdf}
  \caption{Both averages for all three engines. They are shown together
  because they disagree about which classical engine leads and agree about which
  came last, and a chart carrying one of them would pick a winner by picking a
  denominator.}
  \label{fig:engine-f1}
\end{figure}

The language model reached \vLlmMacroF{} macro-F1 and \vLlmMicroF{} micro-F1;
the n-gram model \vNgramMacroF{} and \vNgramMicroF{}; the rule table
\vRulesMacroF{} and \vRulesMicroF{}.\footnote{Every number in this paragraph is a
generated macro. The values are in Table~\ref{tab:headline} and the paths they
came from are printed under it.} On this corpus, at this quantization, with this
prompt, the small linear model is not the compromise. It is the better
classifier, and it is also several thousand times faster, needs no accelerator,
and adds nothing meaningful to an install.

The fourth row of Table~\ref{tab:headline} reads \emph{absent} in words rather
than being left blank. A blank cell in a comparison table reads as a measurement
of zero, which would be a fabricated number in the one place a reader trusts
most.

That row is absent for a different reason from the one it originally carried, and
the difference is worth a sentence here rather than only later. The transformer
rung was specified as an optional later phase, and it was then built, fine tuned,
quantized and measured against this same test split. It is the most accurate
classifier this project has produced and it did not ship, because it missed a
latency budget that was fixed before it was trained. This table is the three
engine benchmark it was not part of. Its own run, its result file, and the
arithmetic that rejected it are in Section~\ref{sec:transformer-rung}.

\section{The two averages disagree, and the disagreement is the mechanism}

Macro-F1 says the rule table wins. Micro-F1 says the n-gram model wins. Both are
in the same files and both are correct.

The mechanism is abstention. The rule table answers \code{UNKNOWN} on
\vRulesAbstention{} of the fields and is right on \vRulesCommittedAccuracy{} of
what it does answer. The n-gram model abstains on \vNgramAbstention{} and is
right on \vNgramCommittedAccuracy{}. Abstention costs almost nothing under macro
averaging, where an abstention is a miss on one class, and costs a full share of
the field weighted denominator under micro averaging. A model that answers nearly
everything therefore earns field weighted accuracy and pays label weighted
average every time it guesses a rare class wrongly.

\begin{table}[htbp]
  \centering
  \caption{Abstention, and accuracy on the fields each engine did commit to.}
  \label{tab:abstention}
  \begingroup\small\setlength{\tabcolsep}{4pt}
\begin{tabular}{lrrrr}
\toprule
\textbf{Engine} & \textbf{Fields} & \textbf{Answered UNKNOWN} & \textbf{Abstention rate} & \textbf{Accuracy when committed} \\
\midrule
\texttt{rules} & 2317 & 607 & 0.2620 & 0.9702 \\  % results: experiments/results/bench/2026-08-26T20-32-48Z_bench_6457ac7/benchmark.json
\texttt{ngram} & 2317 & 84 & 0.0363 & 0.8204 \\  % results: experiments/results/bench/2026-08-26T20-32-48Z_bench_6457ac7/benchmark.json
\texttt{llm} & 2317 & 59 & 0.0255 & 0.7037 \\  % results: experiments/results/bench/2026-08-26T20-32-48Z_bench_6457ac7/benchmark.json
\bottomrule
\end{tabular}
\endgroup
\par\smallskip\noindent{\footnotesize Source: \rpath{experiments/results/bench/2026-08-26T20-32-48Z_bench_6457ac7/benchmark.json}}

\end{table}

That is not the rule table being timid. On this corpus abstention is the better
policy, and the rule table has it by construction, because its default branch is
a tier that never clears a threshold.

\section{The language model abstained least of all}

The system prompt tells the model to answer \code{UNKNOWN} when the evidence is
insufficient and not to guess. It abstained on \vLlmAbstention{} of fields, less
than either classical engine, and was right on \vLlmCommittedAccuracy{} of what
it committed to, less than either
(Table~\ref{tab:abstention}).

\textbf{Instructing a model to decline is not equivalent to giving it a
threshold}, and this is the measurement of that gap. It is also the prediction
that failed in the most useful direction: abstention was registered to land
between the two classical engines and landed below both.

\section{Locale and tier}

\begin{table}[htbp]
  \centering
  \caption{Seen locales against the held out locale. The unseen locale column is
  the multilingual claim's actual test: French forms whose templates are also
  unseen.}
  \label{tab:slices}
  \begingroup\small\setlength{\tabcolsep}{4pt}
\begin{tabular}{lrrrr}
\toprule
\textbf{Engine} & \textbf{Seen macro-F1} & \textbf{Seen micro-F1} & \textbf{Unseen macro-F1} & \textbf{Unseen micro-F1} \\  % results: experiments/results/bench/2026-08-26T20-32-48Z_bench_6457ac7/benchmark.json
\midrule
\texttt{rules} & 0.7991 & 0.7743 & 0.7276 & 0.7340 \\  % results: experiments/results/bench/2026-08-26T20-32-48Z_bench_6457ac7/benchmark.json
\texttt{ngram} & 0.7735 & 0.8532 & 0.5637 & 0.6649 \\  % results: experiments/results/bench/2026-08-26T20-32-48Z_bench_6457ac7/benchmark.json
\texttt{llm} & 0.6373 & 0.6981 & 0.6045 & 0.6809 \\  % results: experiments/results/bench/2026-08-26T20-32-48Z_bench_6457ac7/benchmark.json
\bottomrule
\end{tabular}
\endgroup
\par\smallskip\noindent{\footnotesize Source: \rpath{experiments/results/bench/2026-08-26T20-32-48Z_bench_6457ac7/benchmark.json}}

\end{table}

\begin{table}[htbp]
  \centering
  \caption{Macro-F1 by locale.}
  \label{tab:by-locale}
  \begingroup\small\setlength{\tabcolsep}{4pt}
\begin{tabular}{lrrrr}
\toprule
\textbf{Locale} & \textbf{Fields} & \textbf{rules} & \textbf{ngram} & \textbf{llm} \\
\midrule
de-DE & 376 & 0.7259 & 0.7814 & 0.5906 \\  % results: experiments/results/bench/2026-08-26T20-32-48Z_bench_6457ac7/benchmark.json
en-GB & 392 & 0.7904 & 0.8128 & 0.6491 \\  % results: experiments/results/bench/2026-08-26T20-32-48Z_bench_6457ac7/benchmark.json
en-NG & 377 & 0.7998 & 0.8595 & 0.6643 \\  % results: experiments/results/bench/2026-08-26T20-32-48Z_bench_6457ac7/benchmark.json
en-US & 392 & 0.8033 & 0.8024 & 0.6480 \\  % results: experiments/results/bench/2026-08-26T20-32-48Z_bench_6457ac7/benchmark.json
fr-FR & 376 & 0.7276 & 0.5637 & 0.6045 \\  % results: experiments/results/bench/2026-08-26T20-32-48Z_bench_6457ac7/benchmark.json
ja-JP & 404 & 0.8285 & 0.7599 & 0.6345 \\  % results: experiments/results/bench/2026-08-26T20-32-48Z_bench_6457ac7/benchmark.json
\bottomrule
\end{tabular}
\endgroup
\par\smallskip\noindent{\footnotesize Source: \rpath{experiments/results/bench/2026-08-26T20-32-48Z_bench_6457ac7/benchmark.json}}

\end{table}

\begin{figure}[htbp]
  \centering
  \includegraphics[width=.85\textwidth]{locale_profile.pdf}
  \caption{Macro-F1 by locale, with the held out locale shaded. It is the only
  column in which the language model leads the n-gram model, and it is also a
  column in which the rule table leads them both.}
  \label{fig:locale-profile}
\end{figure}

The held out locale is where broad pretrained world knowledge was supposed to
pay. It is the one slice where the language model leads the n-gram model,
\vLlmUnseenMacroF{} against \vNgramUnseenMacroF{}, and
Chapter~\ref{ch:statistics} shows that this difference does not survive
correction across the family. It is also a slice where the rule table leads both,
at \vRulesUnseenMacroF{}. A vocabulary of localised regular expressions
generalises to an unseen locale better than either learned system does here,
which is an uncomfortable result and is printed as it stands.

\begin{table}[htbp]
  \centering
  \caption{Macro-F1 by markup quality tier.}
  \label{tab:by-tier}
  \begingroup\small\setlength{\tabcolsep}{4pt}
\begin{tabular}{lrrrr}
\toprule
\textbf{Markup tier} & \textbf{Fields} & \textbf{rules} & \textbf{ngram} & \textbf{llm} \\
\midrule
clean & 520 & 0.8677 & 0.9350 & 0.6696 \\  % results: experiments/results/bench/2026-08-26T20-32-48Z_bench_6457ac7/benchmark.json
partial & 519 & 0.8677 & 0.9290 & 0.6522 \\  % results: experiments/results/bench/2026-08-26T20-32-48Z_bench_6457ac7/benchmark.json
mixed & 639 & 0.8005 & 0.7151 & 0.6352 \\  % results: experiments/results/bench/2026-08-26T20-32-48Z_bench_6457ac7/benchmark.json
hostile & 639 & 0.4681 & 0.5409 & 0.5265 \\  % results: experiments/results/bench/2026-08-26T20-32-48Z_bench_6457ac7/benchmark.json
\bottomrule
\end{tabular}
\endgroup
\par\smallskip\noindent{\footnotesize Source: \rpath{experiments/results/bench/2026-08-26T20-32-48Z_bench_6457ac7/benchmark.json}}

\end{table}

\begin{figure}[htbp]
  \centering
  \includegraphics[width=.78\textwidth]{tier_profile.pdf}
  \caption{The tier profile, which is the most informative single picture in the
  benchmark. The two classical engines have strong and opposite slopes; the
  language model is nearly flat and below both until the tiers converge.}
  \label{fig:tier-profile}
\end{figure}

Figure~\ref{fig:tier-profile} is worth reading slowly.

\textbf{The n-gram model is the best engine on clean and partial markup and the
worst classical engine on mixed markup.} It leads on clean at
\vNgramCleanMacroF{} against \vRulesCleanMacroF{}, and trails on mixed at
\vNgramMixedMacroF{} against \vRulesMixedMacroF{}. The mixed tier is the one
built to look like a real page whose sections were written by different teams at
different times. A model that is better on each tier separately and worse on the
tier that mixes them is worth a paragraph rather than a footnote: it suggests the
model is picking up a per page consistency that the mixed tier deliberately
breaks.

\textbf{The hostile tier is where every engine is worst and where the three
converge.} The rule table falls furthest, to \vRulesHostileMacroF{}, because
stripping labels removes most of what a vocabulary of regular expressions has to
work with. The n-gram model reaches \vNgramHostileMacroF{} and the language model
\vLlmHostileMacroF{}. This is the slice where broad pretrained knowledge was
predicted to help most, and none of the differences on it is significant.

The full locale by tier grid, all twenty four cells for all three engines, is
Table~\ref{tab:locale-by-tier} in Appendix~\ref{app:tables}. Every cell clears
the reporting minimum, so nothing in it is suppressed.

\section{What a developer actually experiences}
\label{sec:accusations}

The finding level numbers matter more than the classifier numbers, because they
are what a user of the tool sees. They run through the decision procedure and the
thresholds rather than reporting raw classifier accuracy.

\begin{table}[htbp]
  \centering
  \caption{Finding level precision and recall. Precision is withheld for the two
  classical engines on the second code because neither made enough accusations to
  clear the reporting minimum, and that asymmetry is the finding rather than a
  formatting inconsistency.}
  \label{tab:findings}
  \begingroup\small\setlength{\tabcolsep}{4pt}
\begin{tabular}{llrrrrr}
\toprule
\textbf{Code} & \textbf{Engine} & \textbf{Accusations} & \textbf{Correct} & \textbf{Needed} & \textbf{Precision} & \textbf{Recall} \\
\midrule
\texttt{MISSING\_AUTOCOMPLETE} & \texttt{rules} & 449 & 449 & 790 & 1.0000 & 0.5684 \\  % results: experiments/results/bench/2026-08-26T20-32-48Z_bench_6457ac7/benchmark.json
 & \texttt{ngram} & 334 & 325 & 790 & 0.9731 & 0.4114 \\  % results: experiments/results/bench/2026-08-26T20-32-48Z_bench_6457ac7/benchmark.json
 & \texttt{llm} & 692 & 540 & 790 & 0.7803 & 0.6835 \\  % results: experiments/results/bench/2026-08-26T20-32-48Z_bench_6457ac7/benchmark.json
\midrule
\texttt{WRONG\_AUTOCOMPLETE} & \texttt{rules} & 28 & 28 & 30 & insufficient data & 0.9333 \\  % results: experiments/results/bench/2026-08-26T20-32-48Z_bench_6457ac7/benchmark.json
 & \texttt{ngram} & 23 & 21 & 30 & insufficient data & 0.7000 \\  % results: experiments/results/bench/2026-08-26T20-32-48Z_bench_6457ac7/benchmark.json
 & \texttt{llm} & 125 & 29 & 30 & 0.2320 & 0.9667 \\  % results: experiments/results/bench/2026-08-26T20-32-48Z_bench_6457ac7/benchmark.json
\bottomrule
\end{tabular}
\endgroup
\par\smallskip\noindent{\footnotesize Source: \rpath{experiments/results/bench/2026-08-26T20-32-48Z_bench_6457ac7/benchmark.json}}

\end{table}

\begin{figure}[htbp]
  \centering
  \includegraphics[width=.6\textwidth]{finding_pr.pdf}
  \caption{Where each engine lands on precision against recall for the missing
  declaration finding. \textbf{This compares three threshold policies, not three
  classifiers.} Each engine speaks under its own rule about when to speak, and
  the spread along the horizontal axis is mostly a picture of those rules.}
  \label{fig:finding-pr}
\end{figure}

\textbf{The rule table was right about every field it accused.} It raised
\vRulesMissingAccusations{} missing declaration accusations and
\vRulesMissingWrong{} of them were wrong, at a precision of
\vRulesMissingPrecision{} and a recall of \vRulesMissingRecall{} against
\vMissingNeeded{} fields that needed one.

\textbf{The n-gram model made this project's first wrong accusations.} It raised
\vNgramMissingAccusations{} and was wrong about \vNgramMissingWrong{} of them.
That is the first time in this project that an engine told a developer to add a
token the answer key disagrees with, and it deserves recording rather than
rounding. The threshold policy did not change and the derivation script was run
unmodified; what changed is that the derived boundary now sits on a model whose
development precision at that point was itself below one, so a test split
precision below one is the design working rather than failing. A precision target
is a target, not a guarantee, and a project that reported the target as though it
were the outcome would be reporting a policy as a measurement.

\textbf{The language model found the most real defects and was wrong about
roughly one accusation in five.} It raised \vLlmMissingAccusations{} accusations,
\vLlmMissingWrong{} of them wrong, at a precision of \vLlmMissingPrecision{} and
the highest recall of the three at \vLlmMissingRecall{}. On a developer tool that
means a fifth of the work it hands you is work you should not do.

On the second accusation code the asymmetry is starker.
\vWrongNeeded{} fields needed it, which clears the reporting minimum of
\vReportingMinimum{}, so every recall is reported. Precision is withheld for both
classical engines because neither made enough accusations to be measured, and is
reported for the language model at \vLlmWrongPrecision{} because it made
\vLlmWrongAccusations{}. The engine whose precision can be reported is the one
that accused often enough to be measured, and that is the finding.

None of this is a property of the language model as a classifier. It is a
property of its threshold block, which is zero and zero because a self reported
confidence is forbidden from gating anything, and which was registered before the
run. \textbf{A finding level comparison between these engines is a comparison of
three policies, and that sentence has to travel with the table.}

\section{Latency, and the context that decides whether it matters}

\begin{table}[htbp]
  \centering
  \caption{Per field latency in microseconds, and the share of each run's wall
  clock spent loading and extracting the page.}
  \label{tab:latency}
  \begingroup\small\setlength{\tabcolsep}{4pt}
\begin{tabular}{lrrrr}
\toprule
\textbf{Engine} & \textbf{p50} & \textbf{p95} & \textbf{p99} & \textbf{Load and extract share} \\  % results: experiments/results/bench/2026-08-26T20-32-48Z_bench_6457ac7/benchmark.json
\midrule
\texttt{rules} & 35.2 & 134.0 & 194.7 & 0.9912 \\  % results: experiments/results/bench/2026-08-26T20-32-48Z_bench_6457ac7/benchmark.json
\texttt{ngram} & 165.6 & 261.5 & 355.1 & 0.9874 \\  % results: experiments/results/bench/2026-08-26T20-32-48Z_bench_6457ac7/benchmark.json
\texttt{llm} & 524189.3 & 674727.4 & 774080.6 & 0.0537 \\  % results: experiments/results/bench/2026-08-26T20-32-48Z_bench_6457ac7/benchmark.json
\bottomrule
\end{tabular}
\endgroup
\par\smallskip\noindent{\footnotesize Source: \rpath{experiments/results/bench/2026-08-26T20-32-48Z_bench_6457ac7/benchmark.json}}

\end{table}

\begin{table}[htbp]
  \centering
  \caption{Where each run's wall clock actually went, in seconds.}
  \label{tab:walltime}
  \begingroup\small\setlength{\tabcolsep}{4pt}
\begin{tabular}{lrrrrr}
\toprule
\textbf{Engine} & \textbf{Load} & \textbf{Extract} & \textbf{Classify} & \textbf{Render} & \textbf{Total} \\
\midrule
\texttt{rules} & 131.7 & 4.4 & 0.1 & 0.01 & 137.4 \\  % results: experiments/results/bench/2026-08-26T20-32-48Z_bench_6457ac7/benchmark.json
\texttt{ngram} & 131.8 & 4.4 & 0.5 & 0.02 & 137.9 \\  % results: experiments/results/bench/2026-08-26T20-32-48Z_bench_6457ac7/benchmark.json
\texttt{llm} & 133.0 & 4.9 & 1280.2 & 0.03 & 2565.8 \\  % results: experiments/results/bench/2026-08-26T20-32-48Z_bench_6457ac7/benchmark.json
\bottomrule
\end{tabular}
\endgroup
\par\smallskip\noindent{\footnotesize Source: \rpath{experiments/results/bench/2026-08-26T20-32-48Z_bench_6457ac7/benchmark.json}}

\end{table}

\begin{figure}[htbp]
  \centering
  \includegraphics[width=.95\textwidth]{latency.pdf}
  \caption{Latency percentiles on a log axis, beside the share of wall time the
  browser takes. The log axis is not a stylistic choice: a linear one is
  unreadable across this range.}
  \label{fig:latency}
\end{figure}

The two classical engines run on a CPU with the inference session warmed and its
threads pinned, and both are comfortably inside a millisecond per field:
\vRulesLatencyNinetyFive{} and \vNgramLatencyNinetyFive{} microseconds at the
ninety fifth percentile. The language model runs on a warmed local accelerator
and reaches \vLlmLatencyNinetyFive{}, about \vLatencyRatio{} times the rule
table's figure at the same percentile.

\textbf{The language model's figure is a share, not a per call measurement.} It
classifies a whole page in one request, so its per field number is a batch's
elapsed time divided by the fields in that batch. Batching amortises the round
trip, which makes the figure flattering rather than harsh, and it is not
comparable against a per call measurement. The engine's own description says so
and the run manifest records it.

The row that turns a latency difference into a product difference is the share of
wall time. For the rule engine the browser is \vRulesLoadShare{} of the run and
for the n-gram engine \vNgramLoadShare{}, so the classifier is a rounding error
on a page load and celebrating the difference between their latency columns would
be celebrating nothing. For the language model the browser is \vLlmLoadShare{}
and the classifier is the rest. One is a classifier you do not wait for. The
other is one you do.

In wall clock, over the identical page set: about \vRulesMinutes{} minutes for
the rule engine, \vNgramMinutes{} for the n-gram engine, and \vLlmMinutes{} for
the language model.

\textbf{This is the finding that survives contact with real pages.} The accuracy
gap depends on the corpus being synthetic. The runtime inversion does not.

\section{What the language model run did, and what it did not}

\begin{table}[htbp]
  \centering
  \caption{The language model run, including the failure mode that did not
  occur.}
  \label{tab:llm-run}
  \begingroup\small\setlength{\tabcolsep}{4pt}
\begin{tabular}{lr}
\toprule
\textbf{Quantity} & \textbf{Value} \\
\midrule
Requests & 480 \\  % results: experiments/results/bench/2026-08-26T20-32-48Z_bench_6457ac7/benchmark.json
Attempts & 480 \\  % results: experiments/results/bench/2026-08-26T20-32-48Z_bench_6457ac7/benchmark.json
Requests needing a retry & 0 \\  % results: experiments/results/bench/2026-08-26T20-32-48Z_bench_6457ac7/benchmark.json
Requests still failing after the retry & 0 \\  % results: experiments/results/bench/2026-08-26T20-32-48Z_bench_6457ac7/benchmark.json
Fields scored UNKNOWN by a schema failure & 0 \\  % results: experiments/results/bench/2026-08-26T20-32-48Z_bench_6457ac7/benchmark.json
Smallest batch & 3 \\  % results: experiments/results/bench/2026-08-26T20-32-48Z_bench_6457ac7/benchmark.json
Largest batch & 23 \\  % results: experiments/results/bench/2026-08-26T20-32-48Z_bench_6457ac7/benchmark.json
Median batch & 8 \\  % results: experiments/results/bench/2026-08-26T20-32-48Z_bench_6457ac7/benchmark.json
Fields per request cap & 40 \\  % results: experiments/results/bench/2026-08-26T20-32-48Z_bench_6457ac7/benchmark.json
Prompt tokens & 677632 \\  % results: experiments/results/bench/2026-08-26T20-32-48Z_bench_6457ac7/benchmark.json
Completion tokens & 159126 \\  % results: experiments/results/bench/2026-08-26T20-32-48Z_bench_6457ac7/benchmark.json
\midrule
Monetary cost & null \\
\bottomrule
\end{tabular}
\endgroup
\par\smallskip\noindent{\footnotesize Source: \rpath{experiments/results/bench/2026-08-26T20-32-48Z_bench_6457ac7/benchmark.json}}

\end{table}

\textbf{Schema compliance was perfect, and that was not predicted.} Across
\vLlmRequests{} requests, \vLlmRetried{} needed a retry, \vLlmFailed{} still
failed after one, and \vLlmUnresolved{} fields were scored \code{UNKNOWN} by a
schema failure. Passing the response schema with the request and letting the
server constrain decoding appears to remove the failure mode outright. The retry
ladder is still there and is still tested against recorded transcripts including
deliberately malformed ones; on this run it never fired. A benchmark whose most
interesting failure mode did not occur is a result, and the honest report of it is
the zero rather than a paragraph about robustness.

\textbf{Cost is recorded as null rather than zero.} The calls were local. Zero is
a measurement and the electricity was not free; null is the absence of one. No
price list was retrieved, so no list price estimate is computed, and the price
table document in this repository deliberately contains no prices.

\section{The engine order in the benchmark was a methodological choice}

The engines ran with the language model first, which is not the order the tables
present. The two classical engines take a couple of minutes each of browser time
with no inference call in them, so running the language model last would have
meant several idle minutes before its first call, against a bounded keep alive
policy. A keep alive expiry inside a measured run would have put a cold model
load into the per field latency distribution. Warming immediately before and
running it first avoids measuring the disk.

\section{The two classical engines reproduced their previous run exactly}

Every prediction and every finding, field for field across all \vTestFields{}
fields, and every metric identical to the previous measurement on the same
corpus. Only the latency columns differ, because those are re-measured rather
than recomputed.

That was not the purpose of re-running them. It is the strongest reproducibility
evidence this project has produced and it is worth the sentence.

\section{The transformer that led the table and did not ship}
\label{sec:transformer-rung}

A distilled multilingual encoder was fine tuned on the same training split,
exported to ONNX, quantized to eight bit integers with dynamic activations, and
measured against this same test split through the same runner
\cite{onnxruntime-quantization}. It scores \vBertMacroF{} macro-F1 and
\vBertMicroF{} micro-F1, which makes it the most accurate classifier this project
has measured, and it is not an engine the tool offers.

The condition that rejected it was fixed before it existed, in the first commit
of the phase, ahead of the training dependencies being installed. It is a
conjunction of three bars and the verdict is arithmetic rather than editorial.
The unseen locale had to beat the n-gram model's \vNgramUnseenMacroF{} by at
least \vBertUnseenMargin{}, that is reach \vBertBarUnseen{}, and it reached
\vBertUnseenMacroF{}. The whole split had to stay within \vBertOverallTolerance{}
of the n-gram model's \vNgramMacroF{}, that is not fall below \vBertBarOverall{},
and it reached \vBertMacroF{}. The ninety fifth percentile latency per field had
to sit at or under \vBertBarLatency{} microseconds, and it measured
\vBertLatencyNinetyFive{}.

Two of three is a no ship, and the condition says so in as many words
specifically because two of three with a wide margin on the two is the case where
it would be tempting to reconsider. The margins are wide. The engine beat the n
gram model on every locale and on every tier: the held out locale by
\vBertUnseenGain{}, the hostile tier by \vBertHostileGain{}, the whole split by
\vBertOverallGain{}, and the recall of \code{MISSING\_AUTOCOMPLETE} by
\vBertMissingRecallGain{} at a precision of \vBertMissingPrecision{} against the
n-gram model's \vNgramMissingPrecision{}. Corrected across a family of
\vBertFamilySize{} comparisons over \vBertPairs{} engine pairs, the held out
locale, the hostile tier and the recall gain survive and the whole split
difference does not.

It still does not ship, because the third bar is a budget rather than a target,
and a budget that bends when the numbers are good was never a budget. Per field
the engine costs about \vBertLatencyRatio{} times the n-gram model at the same
percentile on the same pinned single thread session. In whole page terms that
turns classification from a fraction of a percent of the run into
\vBertClassifyShare{} of it, with the browser still \vBertLoadShare{}. So the
cost is not that the tool would feel slow. It is that the classifier's share of
the work would grow by a factor of forty in a tool whose own measurements say the
classifier was never the bottleneck, while adding a tokenizer, a second model
format, and a model file of \vBertIntEightMegabytes{} megabytes that this
repository cannot hold.

The parity contract was written in two parts before either export existed,
because demanding exactness from a quantizer is demanding that it not work. The
first part passed: the FP32 graph agreed with the trained model on the argmax of
every one of the \vBertParityRows{} development rows, with a maximum absolute
logit difference of \vBertParityDelta{} against a ceiling of
\vBertParityTolerance{}. The second part did not. Dynamic INT8 was required to
agree with the FP32 graph on at least \vBertIntEightAgreementBar{} of development
rows and to give up at most \vBertIntEightDropBar{} of development macro-F1; it
agreed on \vBertIntEightAgreement{} and gave up \vBertIntEightDrop{}. Both bounds
are missed, neither by much, and both are recorded rather than adjusted. Under
the pre-registered rule that is a second and independent failure of the latency
bar, which is defined on the quantized engine.

What the specification predicted about this rung was that the advantage would be
concentrated in the held out locale and the hostile tier and would be small or
negative on the clean tier, where the n-gram model already has the label text.
The first two held: both gains are larger than the whole split gain, which is
what concentration has to mean to be checkable. The third did not. The clean tier
gained \vBertCleanGain{}, where the phase had predicted less than
\vPracticalThreshold{}. It is a small contradiction and it points somewhere: on a
clean page the encoder is not seeing anything the linear model cannot see, so the
gap that remains is about how the two generalise from a label they have both been
shown.

The result that survives this corpus being synthetic is not the accuracy number.
It is the shape of the trade. A pretrained multilingual encoder buys most of its
advantage exactly where a bag of character n-grams is weakest, on a locale it was
never trained on and on markup stripped of its labels, and it charges two orders
of magnitude of arithmetic for it.
